\documentclass[12pt,a4paper]{article}

\usepackage[utf8]{inputenc}
\usepackage[T5]{fontenc}
\usepackage[vietnamese]{babel}
\usepackage{lmodern}
\usepackage{geometry}
\usepackage{amsmath}
\usepackage{booktabs}
\usepackage{longtable}
\usepackage{array}
\usepackage{xcolor}
\usepackage{listings}
\usepackage{hyperref}
\usepackage{enumitem}
\usepackage{caption}

\geometry{left=25mm,right=20mm,top=22mm,bottom=22mm}
\hypersetup{
  colorlinks=true,
  linkcolor=blue,
  urlcolor=blue,
  pdftitle={Thiết kế vật lý CPU 4-bit với core utilization 65\%}
}

\definecolor{codebg}{RGB}{245,247,249}
\definecolor{codegreen}{RGB}{20,110,60}
\lstset{
  basicstyle=\ttfamily\small,
  backgroundcolor=\color{codebg},
  frame=single,
  breaklines=true,
  columns=fullflexible,
  keepspaces=true,
  showstringspaces=false,
  commentstyle=\color{codegreen},
  keywordstyle=\color{blue},
  numbers=left,
  numberstyle=\tiny\color{gray},
  numbersep=8pt
}

\newcommand{\code}[1]{\texttt{\detokenize{#1}}}
\newcommand{\um}{$\,\mu\mathrm{m}$}
\newcommand{\umsq}{$\,\mu\mathrm{m}^{2}$}

\title{\textbf{Báo cáo thiết kế vật lý CPU 4-bit}\\
\large OpenLane 2 -- Sky130 -- Core utilization 65\%}
\author{Sinh viên: \rule{6cm}{0.4pt}}
\date{\today}

\begin{document}

\maketitle
\tableofcontents
\newpage

\section{Mục tiêu báo cáo}

Báo cáo trình bày quá trình chuyển thiết kế CPU 4-bit từ RTL SystemVerilog
sang layout GDSII bằng OpenLane 2. Thiết kế sử dụng công nghệ SkyWater
130\,nm, thư viện standard cell \code{sky130_fd_sc_hd}, chu kỳ clock
20\,ns và mục tiêu core utilization bằng 65\%.

Mục tiêu chính gồm:

\begin{itemize}[noitemsep]
  \item kiểm tra chức năng RTL trước khi triển khai vật lý;
  \item tổng hợp RTL thành netlist standard cell;
  \item xây dựng floorplan với utilization 65\%;
  \item thực hiện placement, clock-tree synthesis và routing;
  \item kiểm tra timing, DRC, antenna và LVS;
  \item giải thích lý do lựa chọn từng tham số quan trọng.
\end{itemize}

\section{Lý do chọn thiết kế CPU 4-bit}

CPU 4-bit được chọn vì có đầy đủ những thành phần cơ bản của một hệ thống
số tuần tự nhưng vẫn đủ nhỏ để quan sát rõ từng giai đoạn physical design.
Thiết kế có datapath, ALU, thanh ghi, bộ đếm chương trình, bộ điều khiển
trạng thái và đường clock. Vì vậy, nó phù hợp để minh họa toàn bộ flow
RTL-to-GDSII mà không làm thời gian chạy và việc phân tích trở nên quá phức
tạp.

Các khối RTL chính là:

\begin{itemize}[noitemsep]
  \item \code{cpu_4bit.sv}: top-level, chứa PC, IR, accumulator và các cờ;
  \item \code{alu_4bit.sv}: thực hiện cộng, trừ, logic, dịch, tăng và giảm;
  \item \code{control_unit.sv}: FSM gồm FETCH, EXECUTE và HALT;
  \item program memory đặt bên ngoài CPU và không được tổng hợp vào core.
\end{itemize}

Giao tiếp top-level gồm clock, reset, bus địa chỉ instruction 4-bit, bus dữ
liệu instruction 8-bit và tín hiệu \code{halted}. Việc tách program memory
ra khỏi core giúp phạm vi physical design tập trung vào logic CPU.

\section{Công cụ và công nghệ}

\begin{longtable}{>{\raggedright\arraybackslash}p{0.30\textwidth}
                  >{\raggedright\arraybackslash}p{0.62\textwidth}}
\toprule
\textbf{Thành phần} & \textbf{Lựa chọn và lý do} \\
\midrule
\endhead
OpenLane 2 & Điều phối flow RTL-to-GDSII và lưu state sau từng bước, thuận
tiện cho việc dừng, kiểm tra và chạy tiếp. \\
OpenROAD & Thực hiện floorplan, placement, CTS, routing và phân tích vật lý. \\
Yosys & Tổng hợp RTL và ánh xạ logic sang standard cell. \\
Sky130A & PDK mã nguồn mở, phù hợp cho học tập và nghiên cứu ASIC. \\
\code{sky130_fd_sc_hd} & Thư viện high-density giúp tạo layout nhỏ gọn cho
CPU có diện tích logic nhỏ. \\
KLayout/Magic & Kiểm tra và quan sát layout; hỗ trợ kiểm tra DRC và stream-out. \\
\bottomrule
\end{longtable}

\section{Cấu hình thiết kế}

File \code{config_util65.json} lưu các tham số gốc của thiết kế 65\%.
Trong quy trình thủ công, file này chỉ là tài liệu tham chiếu; các giá trị
tương ứng được truyền bằng Tcl thay vì đưa file cho một OpenLane flow:

\begin{lstlisting}[language={},caption={Cấu hình OpenLane 2}]
{
  "meta": {
    "version": 2,
    "flow": "Classic"
  },

  "PDK": "sky130A",
  "STD_CELL_LIBRARY": "sky130_fd_sc_hd",

  "DESIGN_NAME": "cpu_4bit",
  "VERILOG_FILES": [
    "dir::alu_4bit.sv",
    "dir::control_unit.sv",
    "dir::cpu_4bit.sv"
  ],

  "CLOCK_PORT": "clk",
  "CLOCK_PERIOD": 20.0,

  "FP_CORE_UTIL": 65,
  "FP_PIN_ORDER_CFG": "dir::pin_order.cfg",
  "ERRORS_ON_UNMATCHED_IO": "both"
}
\end{lstlisting}

\subsection{Giải thích các tham số}

\begin{longtable}{>{\raggedright\arraybackslash}p{0.31\textwidth}
                  >{\raggedright\arraybackslash}p{0.61\textwidth}}
\toprule
\textbf{Tham số} & \textbf{Ý nghĩa và lý do lựa chọn} \\
\midrule
\endhead
\code{flow: Classic} & Chọn flow tiêu chuẩn từ RTL đến GDSII, phù hợp với
thiết kế standard-cell không có macro. \\
\code{PDK: sky130A} & Chọn bộ quy tắc và model công nghệ SkyWater 130\,nm. \\
\code{sky130_fd_sc_hd} & Chọn standard-cell library high-density để giảm
diện tích. \\
\code{DESIGN_NAME} & Phải trùng với module top-level \code{cpu_4bit}. \\
\code{VERILOG_FILES} & Chỉ chứa RTL tổng hợp được; testbench và program
memory mô phỏng không được đưa vào synthesis. \\
\code{CLOCK_PORT: clk} & Xác định clock chính cho STA và clock-tree synthesis. \\
\code{CLOCK_PERIOD: 20.0} & Chu kỳ 20\,ns tương đương tần số mục tiêu
$f=1/T=50$\,MHz. Đây là mục tiêu vừa phải đối với CPU nhỏ trên Sky130. \\
\code{FP_CORE_UTIL: 65} & Yêu cầu diện tích cell ban đầu chiếm khoảng 65\%
vùng core khi OpenLane tính kích thước floorplan. \\
\code{FP_PIN_ORDER_CFG} & Điều khiển cạnh và thứ tự I/O để kết nối logic rõ
ràng và routing thuận lợi. \\
\code{ERRORS_ON_UNMATCHED_IO} & Dừng flow nếu danh sách pin không khớp port
RTL, tránh tạo layout thiếu hoặc sai chân. \\
\bottomrule
\end{longtable}

\section{Lý do chọn core utilization 65\%}

Core utilization được hiểu gần đúng là:

\[
U = \frac{A_{\mathrm{cell}}}{A_{\mathrm{core}}}\times 100\%.
\]

Mục tiêu 65\% tạo ra sự cân bằng giữa diện tích và khoảng trống triển khai.
Nếu utilization quá thấp, core lớn hơn cần thiết, làm tăng diện tích die và
có thể làm dây dài hơn. Nếu utilization quá cao, các cell bị ép sát nhau,
không còn đủ không gian cho buffer, clock tree, tối ưu timing và routing.
Điều đó có thể gây congestion và DRC violation.

Giá trị 65\% được chọn vì thiết kế CPU này nhỏ, không chứa macro và kết quả
signoff cho thấy placement và routing vẫn hoàn thành sạch. Đây là giá trị
mục tiêu của bước floorplan, không nhất thiết bằng utilization cuối cùng.
Sau CTS, timing repair và các bước chèn cell, số lượng và diện tích cell có
thể thay đổi.

\section{Pin placement}

Các bus instruction được bố trí theo hướng luồng dữ liệu:

\begin{itemize}[noitemsep]
  \item \code{imem_addr[3:0]} ở cạnh Bắc;
  \item \code{imem_rdata[7:0]} ở cạnh Nam;
  \item \code{clk} và \code{rst_n} ở cạnh Tây;
  \item \code{halted} ở cạnh Đông.
\end{itemize}

Cách bố trí này nhóm các tín hiệu cùng chức năng, làm sơ đồ dễ đọc và giảm
khả năng các bus giao cắt không cần thiết. OpenLane vẫn căn tọa độ pin vào
routing grid hợp lệ.

\section{Quy trình physical design điều khiển bằng Tcl}

Trong phương pháp này, OpenLane không được dùng để chạy từng flow step.
OpenROAD được mở trực tiếp và toàn bộ thao tác physical design được gọi từ
Tcl console. OpenLane chỉ cung cấp Nix shell chứa đúng phiên bản tool và
PDK. Sau mỗi công đoạn, lệnh \code{write_db} tạo checkpoint ODB để có thể
đóng GUI và tiếp tục sau đó.

\subsection{Kiểm tra RTL}

Mô phỏng chức năng phải hoàn thành trước physical design. Physical design
không sửa được lỗi thuật toán hoặc lỗi FSM trong RTL. Kết quả mong đợi của
testbench là:

\begin{lstlisting}[language={}]
TEST PASSED: ACC = 6
\end{lstlisting}

\subsection{Kết quả synthesis đầu vào}

Yosys đã đọc ba file RTL, elaboration thiết kế, tối ưu logic và ánh xạ sang
cell của \code{sky130_fd_sc_hd}. Physical design thủ công không chạy lại
synthesis mà sử dụng trực tiếp gate-level netlist đã tồn tại:

\begin{lstlisting}[language={}]
runs/manual_pd_01/06-yosys-synthesis/cpu_4bit.nl.v
\end{lstlisting}

Sau synthesis cần kiểm tra:

\begin{itemize}[noitemsep]
  \item không có latch ngoài ý muốn;
  \item không còn logic chưa ánh xạ;
  \item top module đúng là \code{cpu_4bit};
  \item netlist không chứa testbench hoặc program ROM;
  \item clock và reset vẫn được nối đúng.
\end{itemize}

\subsection{Mở OpenROAD GUI}

Mở một GUI trống từ thư mục project:

\begin{lstlisting}[language=bash]
cd /home/minhuy/Workspaces/cpu_4bit

nix-shell /home/minhuy/EDA/openlane2/shell.nix \
  --run 'openroad -gui'
\end{lstlisting}

Không dùng \code{openlane --from}, \code{openlane --to} hoặc một flow step
nào. Các lệnh ở những phần tiếp theo được nhập trực tiếp vào Tcl console
của OpenROAD.

\subsection{Đọc PDK, thư viện và netlist synthesis}

Đầu tiên khai báo đường dẫn PDK, sau đó đọc technology LEF, cell LEF,
Liberty và synthesized netlist:

\begin{lstlisting}[language=tcl]
set run_dir runs/manual_tcl_util65

set pdk /home/minhuy/.volare/volare/sky130/versions/\
0fe599b2afb6708d281543108caf8310912f54af/sky130A

read_lef $pdk/libs.ref/sky130_fd_sc_hd/techlef/\
sky130_fd_sc_hd__nom.tlef
read_lef $pdk/libs.ref/sky130_fd_sc_hd/lef/\
sky130_fd_sc_hd.lef
read_lef $pdk/libs.ref/sky130_fd_sc_hd/lef/\
sky130_ef_sc_hd.lef

read_liberty $pdk/libs.ref/sky130_fd_sc_hd/lib/\
sky130_fd_sc_hd__tt_025C_1v80.lib

read_verilog runs/manual_pd_01/06-yosys-synthesis/\
cpu_4bit.nl.v
link_design cpu_4bit

create_clock -name clk -period 20.0 [get_ports clk]
write_db $run_dir/checkpoint_00_synthesis_loaded.odb

sta::redirect_file_begin $run_dir/checkpoint_00_report.rpt
report_clock_properties
report_design_area
sta::redirect_file_end
\end{lstlisting}

Lệnh \code{link_design} tạo design database từ netlist. Lệnh
\code{create_clock} biểu diễn constraint 20\,ns để OpenROAD có thể phân tích
timing trong các bước placement và CTS. Clock chưa được đặt ở chế độ
propagated trước CTS; giai đoạn này dùng ideal clock. Binary OpenROAD đang
sử dụng không có lệnh \code{redirect} hoặc \code{tee}, vì vậy báo cáo được
lưu bằng cặp lệnh \code{sta::redirect_file_begin} và
\code{sta::redirect_file_end}.

\subsection{Khởi tạo floorplan 65\%}

Floorplan xác định die, core và placement rows. Lệnh Tcl sử dụng là:

\begin{lstlisting}[language=tcl]
initialize_floorplan -utilization 65 -aspect_ratio 1.0 -core_space {10.88 10.88 5.52 5.52} -site unithd
make_tracks

write_db $run_dir/checkpoint_01_floorplan.odb
write_def $run_dir/checkpoint_01_floorplan.def
\end{lstlisting}

\code{-aspect_ratio 1.0} yêu cầu core gần hình vuông. Hình dạng này phù hợp
với một CPU nhỏ không chứa macro và không có yêu cầu package đặc biệt.

Thứ tự là \code{bottom top left right}. Placement site \code{unithd} có
chiều cao 2.72\um và chiều rộng 0.46\um. OpenLane dùng bốn lần chiều cao
site cho biên trên/dưới và mười hai lần chiều rộng site cho biên trái/phải:

\[
\begin{aligned}
M_{\mathrm{bottom}}=M_{\mathrm{top}} &= 4\times2.72=10.88\ \mu\mathrm{m},\\
M_{\mathrm{left}}=M_{\mathrm{right}} &= 12\times0.46=5.52\ \mu\mathrm{m}.
\end{aligned}
\]

Các khoảng trống này dành cho I/O, power distribution và routing gần biên.
Việc sử dụng bội số nguyên của placement site cũng giữ core thẳng hàng với
standard-cell grid.

\code{initialize_floorplan} chỉ tạo die, core và placement rows. Lệnh
\code{make_tracks} phải được gọi sau đó để tạo routing tracks từ pitch và
offset trong technology LEF. Nếu thiếu bước này, \code{place_pins} báo lỗi
\code{PPL-0021} vì không tìm thấy horizontal track trên \code{met3}.

\subsection{Đặt I/O pin}

Đặt các pin lên layer ngang \code{met3} và layer dọc \code{met2}:

\begin{lstlisting}[language=tcl]
set_io_pin_constraint -pin_names {imem_addr[3] imem_addr[2] imem_addr[1] imem_addr[0]} -region top:* -group -order
set_io_pin_constraint -pin_names {imem_rdata[7] imem_rdata[6] imem_rdata[5] imem_rdata[4] imem_rdata[3] imem_rdata[2] imem_rdata[1] imem_rdata[0]} -region bottom:* -group -order
set_io_pin_constraint -pin_names {clk rst_n} -region left:* -group -order
set_io_pin_constraint -pin_names {halted} -region right:*
place_pins -hor_layers met3 -ver_layers met2 -corner_avoidance 5 -min_distance_in_tracks

set block [ord::get_db_block]
set pins_without_geometry {}
foreach bterm [$block getBTerms] {
  if {[llength [$bterm getBPins]] == 0} {
    lappend pins_without_geometry [$bterm getName]
  }
}
puts "Pins without geometry: $pins_without_geometry"

write_db $run_dir/checkpoint_02_pins.odb
write_def $run_dir/checkpoint_02_pins.def
\end{lstlisting}

Trong Tcl API, \code{top}, \code{bottom}, \code{left} và \code{right}
tương ứng Bắc, Nam, Tây và Đông. Dấu ngoặc vuông trong tên bus được bảo vệ
bởi Tcl braces. Các lớp met2/met3 được chọn vì chúng nằm trong routing
stack hợp lệ của Sky130 và cho phép pin kết nối thuận lợi với router.

\subsection{Tapcell và well-tie}

Tapcell nối well/substrate về nguồn để tránh well bị floating và giảm nguy
cơ latch-up. Khoảng cách mục tiêu là 13\um:

\begin{lstlisting}[language=tcl]
tapcell -distance 13 -tapcell_master sky130_fd_sc_hd__tapvpwrvgnd_1 -endcap_master sky130_fd_sc_hd__decap_3 -halo_width_x 10 -halo_width_y 10
write_db $run_dir/checkpoint_03_tapcells.odb
write_def $run_dir/checkpoint_03_tapcells.def
\end{lstlisting}

Kết quả kiểm tra thực tế là 30 tapcell, 36 endcap và tổng cộng 206
instances. Ngay sau tapcell, revision OpenROAD được sử dụng đã in
\code{Design area 0} kèm \code{68\% utilization}. Vì hai trường trong cùng dòng
không nhất quán, diện tích movable cell được kiểm tra trực tiếp từ kích
thước master trong OpenDB. Kết quả là 140 movable instances với diện tích
1591.5264\umsq, xác nhận netlist logic không bị mất.

\subsection{Power planning}

Power planning tạo mạng phân phối \code{VPWR} và \code{VGND}. Mục tiêu là
đưa nguồn đến các hàng standard cell bằng đường kim loại có điện trở đủ
thấp. Sau khi tạo PDN cần kiểm tra connectivity và IR drop. Phân tích IR
trong flow này vẫn mang tính ước lượng vì không có file vị trí nguồn ngoài
\code{VSRC_LOC_FILES}.

\begin{lstlisting}[language=tcl]
add_global_connection -net VPWR -inst_pattern .* -pin_pattern {^VPWR$} -power
add_global_connection -net VPWR -inst_pattern .* -pin_pattern {^VPB$} -power
add_global_connection -net VGND -inst_pattern .* -pin_pattern {^VGND$} -ground
add_global_connection -net VGND -inst_pattern .* -pin_pattern {^VNB$} -ground
global_connect

set_voltage_domain -name CORE -power VPWR -ground VGND
define_pdn_grid -name stdcell_grid -starts_with POWER -voltage_domain CORE -pins {met4 met5}

add_pdn_stripe -grid stdcell_grid -layer met1 -width 0.48 -followpins -starts_with POWER
add_pdn_stripe -grid stdcell_grid -layer met4 -width 1.6 -pitch 153.6 -offset 16.32 -spacing 1.7 -starts_with POWER
add_pdn_stripe -grid stdcell_grid -layer met5 -width 1.6 -pitch 153.18 -offset 16.65 -spacing 1.7 -starts_with POWER

add_pdn_connect -grid stdcell_grid -layers {met1 met4}
add_pdn_connect -grid stdcell_grid -layers {met4 met5}

pdngen
check_power_grid -net VPWR -error_file $run_dir/VPWR_grid_errors.rpt
check_power_grid -net VGND -error_file $run_dir/VGND_grid_errors.rpt
write_db $run_dir/checkpoint_04_pdn.odb
write_def $run_dir/checkpoint_04_pdn.def
\end{lstlisting}

Các giá trị width, pitch và offset phải thỏa design rules của PDK. Trước
khi dùng cho tape-out, cần xác nhận lại chúng bằng DRC và power-integrity
analysis; các giá trị trên phù hợp cho thực hành với core nhỏ.

\subsection{Global placement và detailed placement}

Global placement phân bố cell để tối ưu tổng wirelength và congestion.
Detailed placement căn từng cell vào row và site hợp lệ, đồng thời loại bỏ
overlap. Ở utilization 65\%, placement phải dành đủ không gian cho cell
được thêm vào ở các bước sửa timing và CTS.

\begin{lstlisting}[language=tcl]
set_wire_rc -signal -layer met2
set_wire_rc -clock  -layer met3

global_placement -density 0.75 -timing_driven -routability_driven -pad_left 0 -pad_right 0 -init_wirelength_coef 0.25
detailed_placement
optimize_mirroring
check_placement -verbose

write_db $run_dir/checkpoint_05_placement.odb
write_def $run_dir/checkpoint_05_placement.def
\end{lstlisting}

Trong lần chạy đầu tiên chưa có metrics của một implementation trước đó.
Vì vậy, target global placement ban đầu được chọn bằng heuristic:

\[
D_{\mathrm{GPL}} =
U_{\mathrm{floorplan}} + 5P_{\mathrm{cell}} + 10.
\]

Với $U_{\mathrm{floorplan}}=65\%$ và cell padding
$P_{\mathrm{cell}}=0$, target đầu tiên là $75\%$, tức tham số Tcl bằng
0.75. Phần chênh 10 điểm phần trăm là headroom cho thuật toán global
placement. Sau lần chạy đầu, congestion, overflow và timing report mới
được dùng để quyết định giữ hay điều chỉnh density.

\subsection{Clock-tree synthesis}

CTS chèn buffer và xây dựng cây clock nhằm giảm skew, kiểm soát insertion
delay và giới hạn tải trên đường clock. Clock không nên được routing như
một net tín hiệu thông thường vì nó điều khiển đồng thời nhiều flip-flop.

\begin{lstlisting}[language=tcl]
unset_propagated_clock [all_clocks]

clock_tree_synthesis -root_buf sky130_fd_sc_hd__clkbuf_16 -buf_list {sky130_fd_sc_hd__clkbuf_8 sky130_fd_sc_hd__clkbuf_4 sky130_fd_sc_hd__clkbuf_2} -sink_clustering_size 25 -sink_clustering_max_diameter 50 -sink_clustering_enable

set_propagated_clock [all_clocks]
detailed_placement
check_placement -verbose
write_db $run_dir/checkpoint_06_cts.odb
write_def $run_dir/checkpoint_06_cts.def
\end{lstlisting}

\subsection{Routing}

Global routing ước lượng đường đi và congestion. Detailed routing tạo dây
và via thật trên các lớp kim loại theo design rules của Sky130. Router có
thể lặp lại nhiều lần để sửa short, spacing và các vi phạm hình học khác.

\begin{lstlisting}[language=tcl]
set_routing_layers -signal met1-met5
set_routing_layers -clock  met1-met5
set_global_routing_layer_adjustment * 0.30

global_route -congestion_iterations 50 -verbose

detailed_route -bottom_routing_layer met1 -top_routing_layer met5 -output_drc $run_dir/detailed_route_drc.rpt -droute_end_iter 64 -or_seed 42 -verbose 1

write_db $run_dir/checkpoint_07_routed.odb
write_def $run_dir/checkpoint_07_routed.def
\end{lstlisting}

Trong lần chạy đầu không dùng \code{-allow_congestion}. Nếu global router
còn overflow, cần dừng để điều chỉnh floorplan, placement hoặc routing
resources trước khi detailed routing.

\subsection{Signoff}

Các kiểm tra cuối gồm:

\begin{itemize}[noitemsep]
  \item STA setup/hold ở nhiều corner;
  \item DRC bằng router, Magic và KLayout;
  \item antenna check;
  \item LVS so sánh layout với netlist;
  \item XOR check giữa các biểu diễn layout;
  \item power và IR-drop estimation.
\end{itemize}

\begin{lstlisting}[language=tcl]
report_checks -path_delay min_max -format full_clock_expanded
report_worst_slack -max
report_worst_slack -min
report_tns

check_antennas
check_power_grid -net VPWR
check_power_grid -net VGND

write_verilog -include_pwr_gnd $run_dir/cpu_4bit_manual_final.v
write_db $run_dir/cpu_4bit_manual_final.odb
write_def $run_dir/cpu_4bit_manual_final.def
\end{lstlisting}

OpenROAD cung cấp timing, routing và antenna checks. Signoff DRC bằng
Magic/KLayout, LVS và GDS stream-out cần các Tcl/script tương ứng của từng
tool; chúng không được thay thế chỉ bằng việc lưu ODB.

\section{Sự cố LVS trong flow thủ công và cách khắc phục}

\subsection{Hiện tượng và phương pháp khoanh vùng}

Ở lần chạy thủ công đầu tiên, detailed routing hoàn thành và OpenROAD
không báo routing DRC. Tuy nhiên, Magic extraction báo 10 lỗi:

\begin{lstlisting}
Illegal overlap between obsli1 and viali (types do not connect)
\end{lstlisting}

Netgen sau đó cho thấy số lượng device của layout và schematic bằng nhau,
nhưng connectivity không khớp. Vì vậy lỗi không nằm ở synthesis hay số
lượng standard cell mà nằm ở kết nối hình học sau routing.

Các checkpoint được kiểm tra lần lượt bằng cùng script
\code{Magic.SpiceExtraction} của OpenLane:

\begin{longtable}{>{\raggedright\arraybackslash}p{0.52\textwidth}
                  >{\raggedleft\arraybackslash}p{0.30\textwidth}}
\toprule
\textbf{Checkpoint} & \textbf{Magic extraction} \\
\midrule
\endhead
\code{checkpoint_07_cts.def} & 0 lỗi \\
\code{checkpoint_09_detailed_routing.def} & 0 lỗi \\
\code{checkpoint_10_filler.def} & 10 lỗi \\
\bottomrule
\end{longtable}

Kết quả này xác định filler insertion là thời điểm đầu tiên tạo lỗi. Các
cell decap/filler được chèn vào khoảng trống sau khi dây đã được route.
Một số obstruction \code{obsli1} của các cell mới chèn chồng lên
\code{li1/viali} đã tồn tại.

\subsection{Lệnh cũ gây ra lỗi}

Thứ tự cũ là routing trước, filler insertion sau:

\begin{lstlisting}[language=tcl]
read_db $run_dir/checkpoint_07_cts.odb
global_route -congestion_iterations 50 -verbose
detailed_route -bottom_routing_layer met1 -top_routing_layer met5 -output_drc $run_dir/detailed_route_drc.rpt -droute_end_iter 64 -or_seed 42 -verbose 1
write_db $run_dir/checkpoint_09_detailed_routing.odb
write_def $run_dir/checkpoint_09_detailed_routing.def

filler_placement -prefix FILLER_ {sky130_ef_sc_hd__decap_12 sky130_fd_sc_hd__decap_8 sky130_fd_sc_hd__decap_6 sky130_fd_sc_hd__decap_4 sky130_fd_sc_hd__decap_3 sky130_fd_sc_hd__fill*}
write_db $run_dir/checkpoint_10_filler.odb
write_def $run_dir/checkpoint_10_filler.def
\end{lstlisting}

Thử chạy lại \code{detailed_route} trực tiếp trên checkpoint 10 không
thành công vì database đã chứa routing cũ. Detailed router dừng với
\code{DRT-1010 Unsupported non-orthogonal wire}. Do đó không thể sửa bằng
cách gọi lại detailed routing trực tiếp trên checkpoint này.

\subsection{Thứ tự lệnh mới}

Giải pháp là quay về checkpoint CTS chưa có signal routing, chèn filler
trước, rồi chạy global và detailed routing. Khi đó router nhìn thấy
obstruction của toàn bộ decap/filler và tránh chúng ngay từ đầu.

\begin{lstlisting}[language=tcl]
read_db $run_dir/checkpoint_07_cts.odb
filler_placement -prefix FILLER_ {sky130_ef_sc_hd__decap_12 sky130_fd_sc_hd__decap_8 sky130_fd_sc_hd__decap_6 sky130_fd_sc_hd__decap_4 sky130_fd_sc_hd__decap_3 sky130_fd_sc_hd__fill*}
global_connect
set_routing_layers -signal met1-met5
set_routing_layers -clock met1-met5
set_global_routing_layer_adjustment * 0.30
global_route -congestion_iterations 50 -verbose
detailed_route -bottom_routing_layer met1 -top_routing_layer met5 -output_drc $run_dir/detailed_route_with_fillers_drc.rpt -droute_end_iter 64 -or_seed 42 -verbose 1
write_db $run_dir/checkpoint_11_fillers_before_routing.odb
write_def $run_dir/checkpoint_11_fillers_before_routing.def
\end{lstlisting}

Magic extraction trên checkpoint 11 hoàn thành mà không có extraction
error. Các kiểm tra sau khi sửa cũng cho kết quả antenna bằng 0 và toàn bộ
shape trên hai net \code{VPWR}, \code{VGND} đều connected.

\subsection{Sửa powered Verilog dùng cho LVS}

Lệnh cũ chỉ ghi logical Verilog:

\begin{lstlisting}[language=tcl]
write_verilog $run_dir/cpu_4bit_manual_final_fixed.v
\end{lstlisting}

File tạo ra không có hai top-level power ports và không ghi power/ground
pin của các instance. Vì vậy Netgen báo 152 device ở cả hai phía nhưng số
net là $156$ và $760$.

Lệnh đúng cho LVS phải bao gồm power và ground:

\begin{lstlisting}[language=tcl]
global_connect
write_verilog -include_pwr_gnd $run_dir/cpu_4bit_manual_final_fixed.v
\end{lstlisting}

Powered Verilog mới có 317 kết nối \code{.VPWR(...)} tương ứng với 317
instance trong layout. Kết quả LVS cuối cùng là:

\begin{lstlisting}
Circuit 1 contains 152 devices, Circuit 2 contains 152 devices.
Circuit 1 contains 156 nets,    Circuit 2 contains 156 nets.
Final result:
Circuits match uniquely.
\end{lstlisting}

Như vậy, hai thay đổi cần thiết so với flow cũ là chèn filler trước khi
routing và dùng \code{write_verilog -include_pwr_gnd} để tạo schematic
netlist cho LVS.

\section{Signoff và mô phỏng post-layout của layout đã sửa}

\subsection{Stream-out GDS và DRC}

DEF đã sửa được stream-out lại thành
\code{cpu_4bit_manual_final_fixed.gds}. GDS cũ không được tái sử dụng vì
nó không chứa thứ tự filler/routing mới. KLayout chạy deck
\code{sky130A_mr.drc} với các nhóm FEOL, BEOL và off-grid được bật:

\begin{lstlisting}
klayout -b -zz -r $pdk/libs.tech/klayout/drc/sky130A_mr.drc -rd input=$PWD/runs/manual_tcl_util65/cpu_4bit_manual_final_fixed.gds -rd topcell=cpu_4bit -rd report=$PWD/runs/manual_tcl_util65/reports_fixed/drc_violations.klayout.xml -rd feol=true -rd beol=true -rd floating_metal=false -rd offgrid=true -rd seal=true -rd threads=16
\end{lstlisting}

Report database cuối chứa 0 DRC item. Cùng với LVS match uniquely, kết
quả này xác nhận GDS fixed không có vi phạm được hai signoff deck phát
hiện và connectivity layout khớp powered netlist.

\subsection{RC parasitic extraction và post-layout STA}

OpenRCX dùng nominal rule set của Sky130 để trích xuất điện trở và điện
dung ký sinh từ routing:

\begin{lstlisting}[language=tcl]
read_db $run_dir/cpu_4bit_manual_final_fixed.odb
define_process_corner -ext_model_index 0 CURRENT_CORNER
extract_parasitics -ext_model_file $pdk/libs.tech/openlane/rules.openrcx.sky130A.nom.calibre -lef_res
write_spef $run_dir/cpu_4bit_manual_final_fixed.nom.spef
\end{lstlisting}

Sau khi nạp nominal Liberty, clock 20\,ns và SPEF, post-layout STA cho kết
quả:

\begin{longtable}{>{\raggedright\arraybackslash}p{0.52\textwidth}
                  >{\raggedleft\arraybackslash}p{0.30\textwidth}}
\toprule
\textbf{Chỉ tiêu post-layout nominal} & \textbf{Kết quả} \\
\midrule
\endhead
Setup worst slack & $+16.19$\,ns \\
Hold worst slack & $+0.53$\,ns \\
Total negative slack & $0.00$\,ns \\
\bottomrule
\end{longtable}

Cả setup và hold slack đều dương, đồng thời TNS bằng 0, nên timing đạt ở
nominal corner với clock period 20\,ns.

\subsection{Tạo SDF và gate-level simulation}

SDF dùng cho timing back-annotation được tạo từ trạng thái STA đã nạp
SPEF:

\begin{lstlisting}[language=tcl]
write_sdf -corner nom_tt_025C_1v80 -include_typ -digits 4 $run_dir/cpu_4bit_manual_final_fixed.nom.sdf
\end{lstlisting}

Testbench \code{tb_gls.sv} annotate file này vào instance
\code{cpu_4bit_tb_gls.dut} với chế độ \code{TYPICAL}. Xcelium compile
powered netlist, Sky130 UDP primitives, timing models và testbench bằng:

\begin{lstlisting}
xrun -64bit -sv -define USE_POWER_PINS -timescale 1ns/1ps -access +rwc -xmlibdirname runs/manual_tcl_util65/xcelium_postlayout_fixed.d $pdk/libs.ref/sky130_fd_sc_hd/verilog/primitives.v gls_models/sky130_fd_sc_hd_reset_width.v runs/manual_tcl_util65/cpu_4bit_manual_final_fixed.v tb_gls.sv -top cpu_4bit_tb_gls
\end{lstlisting}

Xcelium báo SDF annotation thành công:

\begin{itemize}[noitemsep]
  \item 722/722 path delays được annotate, đạt 100\%;
  \item 114/114 timing checks được annotate, đạt 100\%;
  \item mô phỏng kết thúc tại 250\,ns với thông báo
  \code{POST-LAYOUT GLS PASSED: ACC = 6}.
\end{itemize}

Do đó flow thủ công fixed đã hoàn thành đủ các mục tiêu: layout,
GDS stream-out, DRC, LVS, parasitic extraction, post-layout STA và
post-layout gate-level simulation.

\section{Kết quả tham chiếu của thiết kế 65\%}

Các số liệu dưới đây lấy từ implementation 65\% đã được kiểm chứng tại
\code{runs/manual_pd_util65/final/metrics.json}. Chúng là mốc tham chiếu để
đánh giá quá trình chạy Tcl thủ công. Sau khi hoàn thành manual flow, cần
trích xuất lại metrics vì việc thay đổi seed, placement option hoặc PDN có
thể làm kết quả khác đi.

\begin{longtable}{>{\raggedright\arraybackslash}p{0.48\textwidth}
                  >{\raggedleft\arraybackslash}p{0.36\textwidth}}
\toprule
\textbf{Chỉ tiêu} & \textbf{Kết quả} \\
\midrule
\endhead
Die area & 4311.44\umsq \\
Core area & 2409.81\umsq \\
Final instance utilization & 67.7051\% \\
Số instance cuối & 196 \\
Tổng instance area & 1631.56\umsq \\
Setup worst slack & $+11.859$\,ns \\
Hold worst slack & $+0.115$\,ns \\
Total routed wirelength & 3137\um \\
Maximum net wirelength & 150.69\um \\
Số via & 1180 \\
Total power ước lượng & 0.4218\,mW \\
Worst IR drop ước lượng & 0.166\,mV \\
Final routing DRC violations & 0 \\
Antenna violations & 0 \\
Magic DRC violations & 0 \\
KLayout DRC violations & 0 \\
LVS errors & 0 \\
XOR differences & 0 \\
\bottomrule
\end{longtable}

\section{Phân tích kết quả}

Utilization cuối bằng 67.7051\%, cao hơn mục tiêu floorplan 65\%. Đây không
phải lỗi. Giá trị 65\% được dùng để tính kích thước core tại thời điểm
floorplan, còn flow có thể chèn hoặc thay đổi buffer, clock cells, tie cells
và cell sửa timing ở những bước sau.

Setup slack và hold slack đều dương, cho thấy thiết kế đáp ứng constraint
clock 20\,ns trong các corner được kiểm tra. Hold slack nhỏ hơn setup slack,
nhưng vẫn dương nên không có hold violation.

Kết quả DRC, antenna và LVS đều bằng 0 cho thấy layout hợp lệ về hình học,
không có lỗi antenna được báo và kết nối vật lý khớp với netlist. Worst IR
drop ước lượng rất nhỏ so với nguồn 1.8\,V. Tuy nhiên, kết luận IR-drop cần
được xem là sơ bộ do chưa mô hình hóa chính xác vị trí nguồn package.

\section{Các lưu ý kỹ thuật}

\begin{itemize}
  \item Không đồng nhất target utilization với final utilization.
  \item Không sửa utilization bằng cách thay đổi trực tiếp die sau khi đã
  placement; nên quay lại checkpoint floorplan hoặc tạo một bộ checkpoint
  mới.
  \item GUI vừa dùng để quan sát vừa cung cấp Tcl console. Người thiết kế
  phải chủ động lưu ODB sau từng công đoạn để flow thủ công có thể lặp lại.
  \item Synthesis state không chứa DEF hoặc ODB. Vì vậy OpenROAD chưa có
  layout để hiển thị cho đến khi chạy \code{initialize_floorplan}.
  \item Cần giữ đường dẫn PDK và version PDK nhất quán khi chạy thủ công.
  \item Trước tape-out thực tế cần có nguồn activity chính xác, package
  model và vị trí nguồn để phân tích power integrity đáng tin cậy hơn.
\end{itemize}

\section{Kết luận}

CPU 4-bit đã được triển khai thành công với target core utilization 65\%
trên Sky130. Flow hoàn thành synthesis, floorplan, placement, CTS, routing
và signoff. Timing setup/hold đều đạt, các kiểm tra DRC, antenna và LVS đều
pass. Kết quả cho thấy 65\% là một target khả thi cho phiên bản RTL hiện
tại, đồng thời tạo ra một layout nhỏ gọn mà router vẫn xử lý thành công.

\end{document}
